% ------------------------------------------------------------------
% Section: An eight-vertex neutral square rotor
% Drafted by: Claude (Fable 5), 2026-07-17
% Sources (all claims trace to these; see claude_eight_vertex_rotor_claims.md):
%   problems/23/writeup/WALL_ATTACK_R39_GPTPRO56.md   (secs. 4-5)
%   problems/23/writeup/_claude_r39_8vtx_rotor_gate.py (exhaustive verifier,
%     SHA-256 6d74bcbd1bab12948c5e1a498f62a7185b03743a2b701ec5aeba6f54b01b2aeb)
% Assumes the host preamble of the companion paper (amsthm environments
% theorem/proposition/lemma/corollary/definition/remark, booktabs,
% \bip, \supp, \mc as in shortest_support_obstructions/main.tex).
% ------------------------------------------------------------------

\section{An eight-vertex neutral square rotor}\label{sec:rotor8}

As before, we fix a cut \(V(G)=V_0\mathbin{\dot\cup}V_1\) of a
triangle-free graph \(G\), write \(B\) for the crossing edges and \(M\)
for the monochromatic (\emph{bad}) edges, and for \(e=xy\in M\) with
\(d_B(x,y)=4\) let \(\supp_B(e)\) be the set of crossing edges lying on
shortest \(x\)-\(y\) paths in \(B\).

The obstructions considered so far are static: sets of bad edges whose
support union is too small. This section records a complementary,
dynamic phenomenon on eight vertices. In the example, support Hall
expansion holds with room to spare and every bad edge even has two
geodesics; yet no choice of one geodesic per bad edge yields
vertex-disjoint paths, and the natural local repair walk, which replaces
a colliding geodesic by an alternative geodesic of the same bad edge,
never terminates: the states form a four-cycle and every state carries a
collision. Every statement in this section is an
elementary finite check, proved in full below; all of them have in
addition been verified exhaustively by machine
(Remark~\ref{rem:rotor8-verification}).

\begin{definition}[The rotor graph]\label{def:rotor8}
Let \(R\) be the graph on the eight vertices \(a,b,p,q,x,y,m,v\) with
the ten edges
\[
 ax,\ yb,\ pm,\ vq,\qquad xm,\ my,\ yv,\ vx,\qquad ab,\ pq,
\]
equipped with the cut \(V_0=\{x,y,p,q\}\), \(V_1=\{m,v,a,b\}\).
The first eight edges cross the cut, so
\(B=\{ax,yb,pm,vq,xm,my,yv,vx\}\) and \(M=\{ab,pq\}\).
The four edges \(xm,my,yv,vx\) form a four-cycle of \(B\), the
\emph{square}; the pendant edges attach \(a,b\) to the opposite corners
\(x,y\) and \(p,q\) to the opposite corners \(m,v\).
\end{definition}

\begin{proposition}\label{prop:rotor8-maxcut}
\(R\) is triangle-free, \(B\) is connected, and the cut of
Definition~\ref{def:rotor8} is a maximum cut:
\[
 \mc(R)=8,\qquad \bip(R)=2,
\]
with bad edges exactly \(ab\) and \(pq\).
\end{proposition}

\begin{proof}
\(B\) is bipartite between \(V_0\) and \(V_1\), so a triangle of \(R\)
must use \(ab\) or \(pq\); these two edges are vertex-disjoint, so it
uses exactly one. But \(N(a)\cap N(b)=\{x,b\}\cap\{y,a\}=\emptyset\)
and \(N(p)\cap N(q)=\{m,q\}\cap\{v,p\}=\emptyset\). Connectivity of
\(B\) is clear: every vertex reaches the square by at most one pendant
edge.

For the maximum cut, consider the four pentagons
\[
 C_1=axmyba,\quad C_2=axvyba,\quad C_3=pmxvqp,\quad C_4=pmyvqp.
\]
Every edge of \(R\) lies in exactly two of them: \(ab,ax,yb\) in
\(C_1,C_2\); \(pq,pm,vq\) in \(C_3,C_4\); and the square edges
\(xm\in C_1,C_3\); \(my\in C_1,C_4\); \(yv\in C_2,C_4\);
\(vx\in C_2,C_3\). A cut crosses at most four edges of a five-cycle,
so for every cut, counting crossing edges once per containing pentagon,
\[
 2\cdot|\{\text{crossing edges}\}|
 =\sum_{i=1}^{4}|\{\text{crossing edges on }C_i\}|\le 4\cdot4=16 .
\]
Hence \(\mc(R)\le8\). The displayed cut crosses all eight edges of
\(B\), so \(\mc(R)=8\) and \(\bip(R)=2\).
\end{proof}

\begin{lemma}[Geodesic families]\label{lem:rotor8-geodesics}
In the fixed cut, both bad edges have \(B\)-distance four. Each has
exactly two \(B\)-geodesics,
\[
 A_m=a\,x\,m\,y\,b,\quad A_v=a\,x\,v\,y\,b \qquad\text{for } ab,
\]
\[
 B_x=p\,m\,x\,v\,q,\quad B_y=p\,m\,y\,v\,q \qquad\text{for } pq.
\]
Consequently \(|\supp_B(ab)|=|\supp_B(pq)|=6\) and
\(\supp_B(M)=B\), so the support Hall condition holds strictly for
every nonempty subset of \(M\).
\end{lemma}

\begin{proof}
The vertices \(a,b,p,q\) have unique \(B\)-neighbours \(x,y,m,v\)
respectively. There is no \(a\)-\(b\) path in \(B\) of length at most
three: length two needs a common neighbour, and
\(N_B(a)\cap N_B(b)=\{x\}\cap\{y\}=\emptyset\); length three needs
\(w\in N_B(x)\cap N_B(b)=\{a,m,v\}\cap\{y\}=\emptyset\). A length-four
path must run \(a\,x\,w\,y\,b\) with
\(w\in N_B(x)\cap N_B(y)\setminus\{a,b\}=\{m,v\}\), giving exactly
\(A_m\) and \(A_v\). Symmetrically, an \(p\)-\(q\) geodesic runs
\(p\,m\,w\,v\,q\) with \(w\in N_B(m)\cap N_B(v)\setminus\{p,q\}
=\{x,y\}\), giving \(B_x\) and \(B_y\). The two geodesics of \(ab\)
share the edges \(ax,yb\), so their union has six edges; likewise for
\(pq\) with \(pm,vq\). The two six-edge supports together cover all
eight edges of \(B\).
\end{proof}

\begin{definition}[Selections, detour moves, owners]\label{def:rotor8-selection}
Let \((V_0,V_1)\) be a cut in which every bad edge has \(B\)-distance
four. A \emph{selection} \(\omega\) assigns to each bad edge one of its
\(B\)-geodesics; its \emph{selected support}
\(S_\omega\subseteq B\) is the union of the chosen edge sets. A
\emph{detour move} at a bad edge \(e\) replaces \(\omega(e)\) by a
different geodesic of \(e\), leaving all other choices unchanged. A
vertex lying on the selected geodesics of two distinct bad edges is an
\emph{owner} of \(\omega\); the selection is \emph{disjoint} if it has
no owners, that is, if the selected geodesics are pairwise
vertex-disjoint.
\end{definition}

A disjoint selection upgrades the edge-counting resource measured by
\(\supp_B\) to a packing of vertex-disjoint odd-cycle witnesses, one
per bad edge. A natural way to look for one is local repair: while two
selected geodesics share a vertex, apply a detour move to one of them.
The rotor graph shows that this repair loop need not terminate, even
when every bad edge offers a spare geodesic.

\begin{theorem}[The four-state rotor]\label{thm:rotor8-states}
In \(R\) with the fixed maximum cut there are exactly four selections,
\[
 \omega_{m,x}=\{A_m,B_x\},\quad
 \omega_{m,y}=\{A_m,B_y\},\quad
 \omega_{v,y}=\{A_v,B_y\},\quad
 \omega_{v,x}=\{A_v,B_x\},
\]
and the following hold.
\begin{enumerate}[label=\textup{(\roman*)}]
\item The state graph on the four selections, with an edge for each
detour move, is the four-cycle
\(\omega_{m,x}\,\omega_{m,y}\,\omega_{v,y}\,\omega_{v,x}\).
\item In the state \(\omega_{u,z}\) the two selected geodesics
intersect exactly in \(\{u,z\}\), which spans an edge of the square;
each of the four square edges arises as the owner pair of exactly one
state. In particular, no selection is disjoint.
\item In every state \(|S_\omega|=7\): exactly one crossing edge is
unselected, namely the square edge vertex-disjoint from the owner pair.
\end{enumerate}
\end{theorem}

\begin{proof}
By Lemma~\ref{lem:rotor8-geodesics} each bad edge has exactly two
geodesics, so the selections are the four listed pairs, and a detour
move toggles exactly one of the two coordinates
\((A_m\leftrightarrow A_v\ \text{or}\ B_x\leftrightarrow B_y)\). Two
states are joined by a move iff they differ in exactly one coordinate,
which is the stated four-cycle; this proves (i). Parts (ii) and (iii)
are read off from the vertex and edge sets of the geodesics:
\begin{center}
\begin{tabular}{c@{\qquad}c@{\qquad}c@{\qquad}c}
\toprule
state & selected geodesics & owner pair & unselected \(B\)-edge\\
\midrule
\(\omega_{m,x}\) & \(A_m,B_x\) & \(\{m,x\}\) & \(yv\)\\
\(\omega_{m,y}\) & \(A_m,B_y\) & \(\{m,y\}\) & \(vx\)\\
\(\omega_{v,y}\) & \(A_v,B_y\) & \(\{v,y\}\) & \(xm\)\\
\(\omega_{v,x}\) & \(A_v,B_x\) & \(\{v,x\}\) & \(my\)\\
\bottomrule
\end{tabular}
\end{center}
For instance
\(V(A_m)\cap V(B_x)=\{a,x,m,y,b\}\cap\{p,m,x,v,q\}=\{m,x\}\) and
\(E(A_m)\cup E(B_x)=\{ax,xm,my,yb\}\cup\{pm,mx,xv,vq\}
=B\setminus\{yv\}\); the other three rows are identical checks.
\end{proof}

Thus every state has exactly one colliding pair of geodesics, sharing
two adjacent vertices, and every detour move merely slides this
collision to an adjacent edge of the square. The following automorphism
shows that the four states are genuinely indistinguishable: no
isomorphism-invariant tie-break distinguishes them.

\begin{proposition}[Neutrality]\label{prop:rotor8-neutral}
The permutation \(\sigma=(x\,m\,y\,v)(a\,p\,b\,q)\) is an automorphism
of \(R\) of order four. It exchanges the two shores, hence preserves
the cut, \(B\), and \(M\); it maps geodesics cyclically,
\[
 A_m\longmapsto B_y\longmapsto A_v\longmapsto B_x\longmapsto A_m ,
\]
and it advances the rotor by one step:
\[
 \sigma\,\omega_{m,x}=\omega_{m,y},\quad
 \sigma\,\omega_{m,y}=\omega_{v,y},\quad
 \sigma\,\omega_{v,y}=\omega_{v,x},\quad
 \sigma\,\omega_{v,x}=\omega_{m,x}.
\]
In particular the four states form a single orbit of a cut-preserving
automorphism: the collision pattern rotates around the square and never
shrinks or grows.
\end{proposition}

\begin{proof}
Direct check on the ten edges:
\(\sigma\) maps
\(ax,yb,pm,vq\) to \(pm,vq,yb,ax\);
the square edges \(xm,my,yv,vx\) to \(my,yv,vx,xm\);
and \(ab,pq\) to \(pq,ab\). Also
\(\sigma(V_0)=\{m,v,b,a\}=V_1\). Applying \(\sigma\) vertexwise to the
paths of Lemma~\ref{lem:rotor8-geodesics},
\(\sigma(A_m)=p\,m\,y\,v\,q=B_y\),
\(\sigma(B_y)=b\,y\,v\,x\,a=A_v\) (reversed),
\(\sigma(A_v)=p\,m\,x\,v\,q=B_x\), and
\(\sigma(B_x)=b\,y\,m\,x\,a=A_m\) (reversed). The action on states
follows by applying \(\sigma\) to each coordinate.
\end{proof}

We call \(R\), with its fixed maximum cut, a \emph{neutral square
rotor}: \emph{square} because the owner pair and the unselected
crossing edge travel around a fixed four-cycle of \(B\), and
\emph{neutral} because an order-four cut-preserving automorphism makes
all states of the repair walk equivalent.

\begin{remark}[What the rotor does and does not show]\label{rem:rotor8-scope}
The rotor is not a Hall violation, let alone a support circuit: a
support circuit has support union one edge smaller than its atom set,
whereas here \(|M|=2\), \(|\supp_B(M)|=8\), and even each single
selected support has seven edges. Moreover the unselected crossing edge
of each state has both endpoints on the square, which meets no bad
edge. What the example eliminates is any argument that upgrades support
abundance to a disjoint system of geodesic representatives, or that
assumes collision repair by detour moves terminates: on \(R\) all four
selections carry the same single collision, and the repair walk never
terminates. Triangle-freeness, maximum-cut optimality, and a
spare geodesic at every bad edge do not preclude this. Whether a rotor
can be embedded in a larger triangle-free graph, keeping the cut
maximum, so that the rotating collision couples to a genuine support
deficiency, is open; no such graph is known. Finally,
\(\bip(R)=2<8^2/25\): as with the static obstructions, the example
constrains proof strategies and has no bearing on the truth of the
Erd\H{o}s \(n^2/25\) conjecture.
\end{remark}

\begin{remark}[Verification]\label{rem:rotor8-verification}
All proofs above are complete and elementary. Independently, every
numbered statement of this section is machine-verified by two
exhaustive scripts included in the ancillary archive. The first
(\texttt{\_claude\_r39\_8vtx\_rotor\_gate.py}, SHA-256 prefix
\texttt{6d74bcbd}) checks triangle-freeness, bipartiteness and
connectivity of \(B\); \(\mc(R)=8\) by enumeration of all \(2^8\)
cuts; completeness of the two geodesic families by exhaustive search
over four-edge crossing paths; and, for each of the four states, the
selected support, the unselected edge, the owner pair, and the
one-coordinate structure of the four detour transitions. The second
(\texttt{verify\_rotor8\_automorphism.py}, SHA-256 prefix
\texttt{756cc65a}) checks the automorphism of
Proposition~\ref{prop:rotor8-neutral}: that \(\sigma\) preserves
adjacency on all vertex pairs, has order exactly four, exchanges the
shores and fixes \(B\) and \(M\) setwise, and---after recomputing
\(d_B=4\) by breadth-first search and re-enumerating the geodesic
families and the four selections from scratch---maps the geodesics and
the states in the stated four-cycles, in particular acting
transitively on the states.
\end{remark}
